\section{敏感性分析}

\subsection{质量组合权重扰动}

问题一的组合权重取熵权：CRITIC $=0.5:0.5$。将比值在 $0.4:0.6$ 至
$0.6:0.4$ 间扰动，重新计算七域质量分：域的排序（book $>$ arxiv $>$
c4 $>$ commoncrawl $>$ github $>$ wikipedia $>$ stackexchange）保持
不变，各域分数最大偏移小于 0.02。说明域级质量排序对权重选择稳健，问题三
中 $Q_0=0.4$ 与域质量映射的取值的置信度较高。

\subsection{广义标度律形式选择}

加性与交互两种质量项形式（\eqref{eq:add}--\eqref{eq:int}）的拟合 $R^2$
分别为 0.9716 与 0.9720，差异仅 $4\times10^{-4}$，且两形式的最优参数几乎
一致（$E,A,a,B,b,C,g$ 差异小于 1\%）。这说明 $Q$ 项结构对结论稳健；选取
交互形式（可退化性更强）不改变问题三的定性结论。弹性与等价性数值对形式
选择的敏感性：交互形式下质量弹性 $-0.141$，加性形式下 $-0.144$，相差
2\% 以内。

\subsection{结构性转移阈值对 $g(Q)$ 形式的敏感性}

三种质量成本函数（指数/幂/对数渐进）下，$Q$ 激活阈值分别
$1.12\times10^{18}$、$3.16\times10^{18}$、$3.55\times10^{18}$。阈值均在
$10^{18}$ 量级、相差不超过 3.2 倍，且都远低于 $10^{19}$ 档预算、远高于
$10^{17}$——故三档预算下的``纯规模 $\to$ 质量激活 $\to$ 质量饱和''转换
结论不依赖 $g(Q)$ 形式。阈值差异可解释：指数型在 $Q_0$ 处斜率最小
（$g'(Q_0)=\gamma\lambda e^{\lambda Q_0}$ 相对小），故更早达到``质量边际
收益 $>$ 边际成本''。

\subsection{$L_{\text{ctx}}$ 敏感性}

见问题三表 4：当 $L_{\text{ctx}}$ 由 2048 增至 131072 时，注意力份额由 5\%
升至 64\%，最优参数量由 0.304B 降至 0.140B，损失由 3.25 升至 3.37。注意力--训练份额相等点
与解析临界 30000 一致（两份额均为 0.401）。敏感性结论：上下文长度是预算
分配的重要外生变量，其不确性直接影响最优规模的量级判断。

\subsection{前沿分解的分位数敏感性}

将前沿回归分位数由 0.9 改为 0.8：$b_N=0.312$、$b_T=0.096$，规模贡献
$0.391/0.487=80.3\%$（0.9 分位时 81.4\%）。分位数选择对``规模占主导''
的定性结论无影响，规模贡献始终在 80\% 附近。

\subsection{自助预测区间对噪声的敏感性}

预测区间宽度与前沿外推噪声 $\sigma$ 成正比：$\sigma=0.10$ 时 24 个月基础
情景区间 $[96.6,121.0]$，$\sigma=0.12$ 时为 $[92.5,125.3]$，$\sigma=0.15$
时为 $[86.8,132.6]$。区间随 $\sigma$ 线性展宽，中位数不变，说明预测点估计
稳健，不确定性主要体现为区间宽度，未改变``24 个月基础情景约 110 分''的核心
结论。

\subsection{跨规模配比收缩的稳健性}

跨尺度收缩指数在 1M/60M/1B 三点拟合时为 $\kappa=0.314$；改用全部 6 个规模点
（含 10b/70b 估计值）拟合得 $\kappa=0.20$；对三点做留一删除，$\kappa$ 在
0.15--0.58 之间波动。量级与符号（配比系数随规模幂律衰减）在不同口径下
稳定成立，但数值精度受规模点选择影响，因此在正文中仅用 $\kappa=0.314$ 作为
方向性结论，不用于定量外推。